\clearpage
\VerdierPhysicalPageAnchor{0250}%
\begingroup\small\setlength{\parskip}{0pt}\setlength{\emergencystretch}{1.5em}
\setlength{\abovedisplayskip}{4pt plus 2pt minus 2pt}%
\setlength{\belowdisplayskip}{4pt plus 2pt minus 2pt}%

\VerdierSemanticAnchor{ch4d-4-6-9-3-isomorphism}
\noindent 例えば、複体 $Y$ が $\operatorname{comp}^{+}(\mathcal A)$ の対象で
関手 $F$ が右定常であるとき、この射は同型である。また、$Y$ について仮定を
置かず、関手 $F$ が定常であるときにも、この射は同型である (4.5.1)。

\VerdierSemanticAnchor{ch4d-4-6-9-derived-functors}
二重複体 $I^{\bullet\bullet}$ は $Y$ の Cartan--Eilenberg 入射分解なので、
複体 $H_1^r(I^{\bullet\bullet})$ は対象 $H^r(Y)$ の入射分解である。
したがって同型

\VerdierSemanticAnchor{ch4d-4-6-9-4}
\[
H_{\mathcal B}^{p}\left(
F^{q-r,\bullet}(H_1^r(I^{\bullet\bullet}))
\right)
\xrightarrow{\sim}
R^pF^{q-r}(H^r(Y)),
\tag{4.6.9.4}
\]
がある。ここで $R^pF^{q-r}\colon\mathcal A\longrightarrow\mathcal B$ は、
関手 $F^{q-r}$ を圏 $\mathcal A$ に制限した関手の $p$ 次右導来関手
（[1] の意味で）を表す。したがって、射 (4.6.9.3) が同型であるとき、同型

\VerdierSemanticAnchor{ch4d-4-6-9-5}
\[
(\mathrm{C.E.})_2^{p,q}\xrightarrow{\sim}
\prod_{r\in\mathbb Z}R^pF^{q-r}(H^r(Y)).
\tag{4.6.9.5}
\]
がある。

\VerdierSemanticAnchor{ch4d-4-6-10}
\medskip\noindent\textbf{4.6.10.}\quad
圏 $\mathcal B$ が適当な宇宙に属するアーベル群の圏 $\mathbf{Ab}$ であり、
関手

\[
F\colon D(\mathcal A)\longrightarrow\mathbf{Ab}
\]
が関手

\VerdierSemanticAnchor{ch4d-4-6-10-ext0}
\[
Y\longmapsto F(Y)=\operatorname{Ext}^{0}(Z,Y),
\]
である場合を考える。ここで $Z$ は $D(\mathcal A)$ の対象である。
$Y$ を $\mathcal A$ の複体とし、
$\varepsilon\colon Y\longrightarrow I^{\bullet\bullet}$ を
Cartan--Eilenberg 入射分解とする。このとき射 (4.6.9.3) は同型であり
(4.6.3)、したがって同型

\VerdierSemanticAnchor{ch4d-4-6-10-1}
\[
(\mathrm{C.E.})_2^{p,q}\xrightarrow{\sim}
\prod_{r\in\mathbb Z}
\operatorname{Ext}^{p}\bigl(H^{r-q}(Z),H^r(Y)\bigr).
\tag{4.6.10.1}
\]
がある。

\VerdierSemanticAnchor{ch4d-4-6-10-stationarity}
また、$F$ に関する Cartan--Eilenberg スペクトル系列は、$Y$ が
$\operatorname{comp}^{+}(\mathcal A)$ の対象で $Z$ が
$D^{-}(\mathcal A)$ の対象であるとき、または Cartan--Eilenberg 入射分解が
有限長であるとき、定常である。さらに、これらの仮定のもとで、この
スペクトル系列の収束先は (4.6.8) により

\VerdierSemanticAnchor{ch4d-4-6-10-abutment}
\[
\prod_{r\in\mathbb Z}
\operatorname{Ext}^{p}\bigl(H^{r-q}(Z),H^r(Y)\bigr)
\Longrightarrow
\operatorname{Ext}^{p+q}(Z,Y).
\]
\endgroup
